\bookchapter{الفصل 32: دليل الإسعافات الأولية لأمان الهاتف}

\chapterintro{هذا الفصل هو «خطة طوارئ». أتمنى ألا تضطر لاستخدامه أبدًا، لكن عند فقدان الهاتف أو سرقته أو التعرض للاحتيال، سيساعدك هذا الدليل على تقليل الخسائر إلى أدنى حد؛ وعند وقوع الزلازل والأمطار الغزيرة، يكون الهاتف أداة لاستقبال التحذيرات وتحديد الموقع واستمرار التواصل. ننصح بقراءته الآن، ثم حفظ لقطات للخطوات الأساسية أو طباعتها ووضعها في المنزل — لأن هاتفك قد لا يكون معك عند وقوع الحادث، وقد يكون من دون بطارية أو إشارة أثناء الكوارث.}

\booksection{السيناريو صفر: الزلازل/الطقس القاسي — كيف يساعدك الهاتف على النجاة}

الهدف هو سلامة الأشخاص وليس أموال الحسابات. لتفاصيل إعدادات التحذير راجع الفصل 13 «تحذيرات الطقس والكوارث الطبيعية والزلازل»؛ هنا نركز على ما تفعله بعد صدور الإنذار.

\textbf{عند صدور إنذار الزلزال: احتمِ أولًا ثم انسحب، ولا تكتفِ بتصوير الفيديو}\par

\begin{enumerate}
  \item النجاة القريبة: انبطح، احتمِ، أمسك بقوة (تحت الطاولة، في زاوية الجدار الحامل، بجانب أثاث متين)، واحمِ رأسك ورقبتك؛ ابتعد عن الزجاج والجدران الخارجية والثريات.
  \item لا تندفع على الدرج بشكل عشوائي: الهزة الرئيسية غالبًا بضع عشرات من الثواني، والدرج قد يكون أكثر خطورة؛ والمصاعد ممنوعة تمامًا.
  \item بعد توقف الهزة انسحب بترتيب: عبر الدرج، وانتبه للهزات الارتدادية والأشياء المتساقطة؛ واتجه إلى مكان مفتوح وابتعد عن أعمدة الكهرباء واللوحات والمباني المتهدمة.
  \item استخدم الهاتف بحكمة: إرسال رسالة قصيرة لجهات الاتصال الطارئة («أنا في XX، بخير/مصاب، أحتاج مساعدة») أفضل من مكالمة طويلة؛ وفضّل النصوص عند توفر الرسائل/WeChat.
  \item عند عدم توفر الإشارة: راجع الفصل 7 «SOS عبر الأقمار الصناعية / الرسائل الفضائية» (إذا كان الطراز يدعمها)؛ أو اصعد إلى مكان مرتفع/مفتوح وحاول مرة أخرى؛ وفي بعض السيناريوهات يساعد جهاز الاتصال اللاسلكي عبر Bluetooth دون شبكة.
  \item أوقف الوظائف غير الضرورية: فعّل وضع الطيران ثم شغّله بشكل دوري، وأوقف مسح Bluetooth، وخفض السطوع، لتدوم البطارية ساعات إضافية؛ واترك آخر شحنة لرسائل الاطمئنان والإنقاذ.
\end{enumerate}

\begin{tcbwarn}{الأمطار الغزيرة/الإعاصير/السيول: تحرك في مرحلة التحذير}
\begin{itemize}
  \item عند صدور تحذير أحمر: قلل الخروج قدر الإمكان، وابتعد عن الأماكن تحت الأرض والجسور المنخفضة والمجاري والمنحدرات.
  \item إذا علقت في الخارج: اتجه إلى أعلى المباني المتينة، ولا تعبر المياه قسرًا (قد تختفي أغطية البئر وقد توجد تسريبات كهربائية)؛ وعند الاتصال بـ 119/110 اذكر المعالم ومستوى المياه.
  \item قبل السفر: استخدم تطبيق الطقس وحسابات التحذير المحلية لتأكيد المخاطر؛ وتجنب في الرحلات الطويلة بالسيارة الطرق الجبلية القريبة من المياه والمنحدرات.
  \item الإبلاغ المتبادل بين أفراد العائلة: اتفقوا على «إذا انقطع التواصل، نعتمد على آخر موقع في WeChat/الرسائل»، لتجنب الاتصال المحموم واستنزاف البطارية.
\end{itemize}

\end{tcbwarn}

\begin{tcbpractice}{4 أمور تفعلها مسبقًا}
\begin{enumerate}
  \item فعّل تحذير الزلازل في النظام / تطبيق تحذير موثوق، واسمح له بكسر وضع الصمت.
  \item ثبّت 2-3 «جهات اتصال طارئة» في أعلى جهات الاتصال، ويمكن عرضها على شاشة القفل (بطاقة الطوارئ الطبية في iOS/Android أو معلومات شاشة القفل).
  \item احتفظ ببطارية متنقلة في المنزل والحقيبة؛ وقبل السفر الطويل اشحنها بالكامل وأحضر الكابل.
  \item احفظ رقمًا واحدًا على الأقل «يمكن الاتصال به دون جهات الاتصال» (رقم جوال أحد أفراد العائلة)، واكتبه على ورق وضعه في المحفظة.
\end{enumerate}

\end{tcbpractice}

\booksection{السيناريو 1: فقدان الهاتف/سرقته — الـ 30 دقيقة الذهبية}

أول 30 دقيقة بعد فقدان الهاتف هي «الوقت الذهبي»، والتعامل الصحيح يحافظ على أموالك وخصوصيتك. اتبع الترتيب التالي ولا تخطِ أي خطوة:

\textbf{الإجراءات القياسية خلال 30 دقيقة بعد الفقد}\par

\begin{enumerate}
  \item اتصل برقمك من هاتف شخص آخر (الدقيقة 1): قد يكون الهاتف سقط قريبًا، وقد يجده شخص طيب ويعيده.
  \item قفل الهاتف عن بُعد (الدقيقتان 2-5): استخدم هاتفًا أو جهاز كمبيوتر آخر لتسجيل الدخول إلى خدمة البحث: iPhone: iCloud.com > ابحث > اختر جهازك > علّم كمفقود
  \item علامات Android: Huawei cloud.huawei.com / Xiaomi i.mi.com / OPPO cloud.oppo.com وغيرها
  \item إيقاف بطاقة SIM (الدقائق 5-10): اتصل بخدمة العملاء (China Mobile 10086 / China Unicom 10010 / China Telecom 10000)، وأبلغ بفقدان الهاتف واطلب إيقاف خدمة SIM. هذه الخطوة هي الأهم — فالمحتال الذي يحصل على بطاقتك يستطيع استقبال رموز التحقق وإعادة تعيين جميع كلمات مرورك.
  \item تجميد حسابات الدفع (الدقائق 10-20): Alipay: اتصل بـ 95188 للإبلاغ، أو سجّل الدخول على هاتف آخر > أنا > الإعدادات > مركز الأمان > تجميد طارئ
  \item WeChat: بعد الاتصال بـ 110 تواصل مع خدمة عملاء WeChat، أو من مركز أمان WeChat > تجميد الحساب
  \item تطبيق البنك: اتصل بخدمة العملاء المطبوعة على ظهر البطاقة البنكية، واطلب تجميد الخدمات المصرفية عبر الإنترنت
  \item إبلاغ الأهل والأصدقاء (الدقائق 20-25): انشر في Moments أو اطلب من العائلة إرسال رسالة لإبلاغهم بفقدان الهاتف، لمنع المحتال من انتحال شخصيتك والاحتيال على الآخرين.
  \item الإبلاغ للشرطة (الدقائق 25-30): اتصل بـ 110 أو اذهب إلى أقرب قسم شرطة لتحرير محضر، واحصل على إيصال التقرير. ستحتاج إليه لاحقًا لمطالبات التأمين ونزاعات البنوك.
\end{enumerate}

\booksection{السيناريو 2: التعرض للاحتيال — الفترة الذهبية 72 ساعة}

بعد اكتشاف أنك تعرضت للاحتيال، الأهم هو الهدوء. القرارات المزيفة أثناء الذعر تسبب خسائر ثانوية. اتبع الخطوات التالية:

\begin{enumerate}
  \item أوقف فورًا جميع عمليات التحويل: مهما قال المحتال، لا تحوّل أي مبلغ إضافي. «حوّل ضمانًا إضافيًا واسترد أموالك السابقة» هو أسلوب احتيال ثانٍ كلاسيكي.
  \item احتفظ بجميع الأدلة: لا تحذف سجل المحادثات! ولا تحذف سجل المكالمات! التقط لقطات شاشة للمحادثات وسجلات التحويل ومعلومات حساب الطرف الآخر. كلها أدلة حاسمة للتبليغ واسترداد الأموال.
  \item اتصل فورًا بـ 96110: هذا هو الخط الساخن الوطني لمكافحة الاحتيال. سيرشدك المشغل إلى إجراء وقف الدفع الطارئ — إذا كان التحويل حديثًا (عادةً خلال 24 ساعة)، فقد يتمكن البنك من تجميد حساب الطرف الآخر وإيقاف الأموال.
  \item اتصل بـ 110 للتبليغ: أحضر الهاتف وبطاقة الهوية وإثبات التحويل إلى أقرب قسم شرطة لتحرير المحضر.
  \item تواصل مع البنك/منصة الدفع: أبلغ بحالة الاحتيال واطلب تجميد الحسابات المشمولة. وإذا كانت بطاقة الائتمان قد سُرقت، اطلب من البنك بدء عملية الرفض.
  \item غيّر جميع كلمات المرور: على جهاز آمن آخر غيّر كلمة مرور Apple ID/حساب العلامة التجارية وWeChat وAlipay وتطبيقات البنوك، وفعّل التحقق الثنائي.
\end{enumerate}

\booksection{السيناريو 3: إصابة الهاتف بفيروس/حصان طروادة — ماذا تفعل؟}

إذا ظهرت على هاتفك أعراض مثل الإعلانات المنبثقة الغامضة أو تثبيت تطبيقات غريبة تلقائيًا أو ارتفاع فاتورة الهاتف أو استهلاك بيانات غير طبيعي، فمن المحتمل أن يكون مصابًا ببرمجيات خبيثة:

\begin{itemize}
  \item الخطوة الأولى: ادخل الوضع الآمن — في Android اضغط مطولًا على زر التشغيل للإيقاف، ثم اضغط مطولًا للتشغيل، وبعد الاهتزاز اضغط مطولًا على زر خفض الصوت حتى تدخل الوضع الآمن (تظهر «الوضع الآمن» في الزاوية اليسرى السفلية). في هذا الوضع يتم تعطيل جميع تطبيقات الطرف الثالث، مما يسهل التشخيص.
  \item الخطوة الثانية: احذف التطبيقات المشبوهة — في الوضع الآمن ادخل الإعدادات > إدارة التطبيقات، ورتّب حسب الوقت، وابحث عن التطبيقات غير المعروفة المثبتة مؤخرًا واحذفها.
  \item الخطوة الثالثة: ثبّت تطبيق أمان وافحص — نزّل «Tencent Mobile Manager» أو «360 Mobile Guard» وقم بفحص كامل.
  \item الخطوة الرابعة: إعادة ضبط المصنع (الحل النهائي) — إذا لم تنجح الطرق السابقة، فخذ نسخة احتياطية من البيانات المهمة ثم أعد الضبط. ملاحظة: إعادة الضبط تحذف جميع البيانات، لذا انسخ الصور وجهات الاتصال والملفات المهمة أولًا.
  \item iPhone نادرًا ما يصاب بالفيروسات — آلية الحماية في iOS تجعل بقاء البرمجيات الخبيثة صعبًا. إذا ظهر خلل في iPhone، فمن المرجح أن يكون خطأً في تطبيق أو إعدادات، وليس فيروسًا. إعادة ضبط المصنع تحل 99\% من المشكلات.
\end{itemize}

\booksection{السيناريو 4: سرقة الحساب — استعد السيطرة بسرعة}

\begin{enumerate}
  \item غيّر كلمة المرور فورًا: إذا كنت لا تزال قادرًا على تسجيل الدخول، فغيّرها في أسرع وقت. كلمة المرور الجديدة يجب أن تكون مختلفة تمامًا عن القديمة.
  \item اطرد الأجهزة الأخرى: في إعدادات أمان الحساب ابحث عن «إدارة أجهزة الدخول»، واطرد جميع الأجهزة غير المعروفة.
  \item افحص سجل تسجيل الدخول: راجع سجلات الدخول خلال آخر 7 أيام، وإذا وجدت دخولاً من موقع آخر، فهذا يعني أن الحساب قد سُرق فعلاً.
  \item فعّل التحقق الثنائي: إذا لم تكن قد فعلت ذلك، ففعّله الآن فورًا. هذه هي الطريقة الأكثر فعالية لمنع سرقة الحساب مرة أخرى.
  \item افحص معلومات الربط: تحقق مما إذا كان رقم الهاتف أو البريد الإلكتروني المرتبط بالحساب قد تغير. وإذا تغير، استخدم وظيفة استرداد الحساب.
  \item افحص إعدادات الدفع: إذا كان الحساب المسرق هو WeChat/Alipay، فتحقق مما إذا كانت المدفوعات بدون كلمة مرور أو الخصم التلقائي قد عُدّلت، وإذا تمت إضافة بطاقات بنكية إلى حسابات غريبة.
\end{enumerate}

\booksection{السيناريو 5: استلام رسالة ابتزاز — لا تخف ولا تدفع}

هناك نوعان شائعان من الابتزاز: (1) «لديّ فيديو/صور مخلة لك، إن لم تدفع سأرسلها لجميع جهات اتصالك»؛ (2) «هاتفك تحت سيطرة قراصنة، إن لم تدفع سأقفله». في هذه الحالة:

\begin{itemize}
  \item 99.9\% احتيال: ليس لديهم أي بيانات خاصة بك فعلًا. يرسلون نفس رسالة الابتزاز لعشرات الآلاف، ودائمًا من يخاف ويدفع.
  \item لا ترد ولا تدفع: الرد يعني إخبار الطرف الآخر بأن «هذا الرقم مستخدم»، وستصبح «هدف ابتزاز» مع مزيد من الإزعاج لاحقًا.
  \item التقط لقطة شاشة كدليل: ثم احظر الطرف الآخر.
  \item غيّر كلمة المرور: إذا كنت قلقًا، فغيّر كلمة مرور حسابك الرئيسي للراحة النفسية.
  \item إذا كان هناك خطر حقيقي لتسريب الخصوصية: اتصل بالشرطة. لديهم خبرة في التعامل مع هذه القضايا.
\end{itemize}

\booksection{السيناريو 6: احتيال تبديل الوجه/تقليد الصوت بالذكاء الاصطناعي — الرؤية لم تعد حقيقة}

منذ 2024، أصبحت تقنيات تبديل الوجه وتقليد الصوت بالذكاء الاصطناعي قادرة على «تحويل» المحتال إلى وجه وصوت عائلتك أو مديرك أو صديقك في مكالمة الفيديو. «يا أمي، وقعت في مشكلة، حوّلي المال بسرعة» — الشخص في الفيديو يبدو ويبدو وكأنه ابنك تمامًا، لكنه مزيف. هذه واحدة من أكثر أساليب الاحتيال صعوبة وخطورة حاليًا.

\textbf{الخطوة الأولى: كشف الثغرات — أربع ثغرات شائعة في التزوير بالذكاء الاصطناعي}\par

\begin{itemize}
  \item معدل رمش العين غير طبيعي: يرمش الإنسان الحقيقي 15-20 مرة في الدقيقة، بينما تبديل الوجه بالذكاء الاصطناعي غالبًا أقل رمشًا أو بإيقاع متيبس.
  \item عدم تزامن حركة الشفاه مع الصوت: بين حركة الشفاه والصوت المسموع يوجد انزياح بجزء من الثانية.
  \item ضبابية حواف الوجه: محيط الوجه وقرب الأذن وأطراف الشعر يظهر فيها «حواف خشنة» أو وميض.
  \item عدم تطابق اتجاه الإضاءة: الإضاءة والظلال على الوجه لا تتوافق مع اتجاه الضوء في الخلفية، أو تظهر الملامح «تتبع» بشكل خاطئ عند الالتفاف.
\end{itemize}

\begin{tcbinfo}{الخطوة الثانية: ثلاث طرق للتحقق من صحة الطرف الآخر}
\begin{itemize}
  \item اطلب حركة معينة: اطلب من الطرف التحول، أو تغطية الوجه، أو وضع اليد على الفم ثم فتحه وإغلاقه. تبديل الوجه بالذكاء الاصطناعي ينكشف بسهولة عند التغطية والدوران بزاوية كبيرة — فالوجه يتشوه أو لا يتبع الحركة.
  \item اسأل أسئلة لا يعرفها إلا الطرف الآخر: اخترع سؤالًا، مثل «أي مستشفى أخذنا فيه الكلب المرة الماضية؟» — إذا كان محتالًا، فلن يجيب أو سيتلعثم.
  \item اتصل بالرقم المعروف: أغلق مكالمة الفيديو، واتصل بالرقم المحفوظ في جهات اتصالك مباشرة. المحتال لا يستطيع التحكم في رقم الهاتف الحقيقي، فينكشف فورًا.
\end{itemize}

\end{tcbinfo}

\begin{tcbwarn}{الخطوة الثالثة: إجراءات التبليغ بعد تأكيد الاحتيال}
\begin{enumerate}
  \item أوقف التحويل فورًا: حتى لو كنت «مستعجلًا»، توقف أولًا. كل من يحثك على «التحويل السريع وعدم قطع المكالمة» هو في الغالب محتال.
  \item احتفظ بجميع الأدلة: سجّل الشاشة والتقط لقطات لشاشة مكالمة الفيديو وسجل المحادثات وحساب الطرف الآخر وإثبات التحويل — هذه مواد حاسمة للخبرة الفنية.
  \item اتصل بـ 96110 الخط الساخن لمكافحة الاحتيال: أبلغ بأنه احتيال «تبديل وجه/تقليد صوت بالذكاء الاصطناعي»، وسيرسل المشغل إجراء وقف الدفع الطارئ.
  \item اتصل بـ 110 للتبليغ: أحضر جميع الأدلة إلى قسم الشرطة لتحرير المحضر، واذكر بوضوح استخدام تقنية التزوير بالذكاء الاصطناعي، لتسهيل فتح القضية وتصنيفها.
  \item أبلغ الأهل والأصدقاء: قد يستخدم المحتال الفيديو المزيف لخداع الآخرين، لذا انشر تحذيرًا مسبقًا في Moments أو مجموعة العائلة.
\end{enumerate}

\end{tcbwarn}

\booksection{السيناريو 7: هجوم SIM swap (خطف SIM) — «سرقة» رقمك}

هذا نوع من الهجوم لم يسمع به كثيرون لكنه خطير جدًا. SIM swap (خطف بطاقة SIM) يعني أن المهاجم يزيف معلومات هويتك وينتحلك للحصول على بطاقة SIM جديدة من مكتب الخدمة. بمجرد نجاح العملية، تختفي إشارة هاتفك فجأة، ويصبح هاتف المهاجم بالبطاقة الجديدة قادرًا على استقبال جميع رموز التحقق — Alipay والبنوك وWeChat، كلها مفتوحة.

\begin{tcbinfo}{كيف تكتشف؟ ثلاث إشارات تحذير لخطف SIM}
\begin{itemize}
  \item يظهر الهاتف فجأة «لا توجد خدمة»/«لا توجد بطاقة SIM»، رغم أنك لم تلمس البطاقة، وإعادة التشغيل لا تجدي — هذا غالبًا لحظة الخطف.
  \item استلام رسالة من شركة الاتصالات تفيد بنجاح «استبدال البطاقة» أو «تنفيذ الخدمة» — رغم أنك لم تطلب ذلك.
  \item توقف مفاجئ عن استقبال رموز التحقق، مع توقف استقبال الرسائل الأخرى — لأن الرسائل قد اختُطفت إلى بطاقة المهاجم.
\end{itemize}

\end{tcbinfo}

\begin{tcbwarn}{إجراءات الحماية — «أقفل» بطاقة SIM أيضًا}
\begin{itemize}
  \item اضبط رمز PIN لبطاقة SIM: المسار الإعدادات > الهاتف > رمز PIN لبطاقة SIM. بعد التفعيل، يجب إدخال PIN عند كل إعادة تشغيل أو تغيير بطاقة. حتى لو حصل المهاجم على البطاقة فلا يستطيع استخدامها.
  \item فعّل التحقق الثنائي لدى شركة الاتصالات: سجّل الدخول إلى تطبيق China Mobile/Unicom/Telecom، وابحث عن «مركز الأمان» أو «التحقق من الخدمات»، وفعّل تنبيهات استبدال البطاقة والتحقق الثنائي للخدمات. عندها سيطلب موظف المكتب تحققًا إضافيًا من الهوية عند استبدال بطاقتك.
  \item الحسابات المهمة لا تعتمد على الرسائل النصية فقط: للحسابات عالية القيمة مثل البنوك وApple ID، فضّل تطبيقات المصادقة (Google Authenticator أو Microsoft Authenticator) أو المفاتيح المادية، بدلًا من رموز التحقق النصية.
  \item أضف علامة مائية إلى نسخ الهوية: اكتب على النسخة المطلوبة «للاستخدام في خدمة XX فقط» — فالمهاجم غالبًا ينتحل هويتك باستخدام معلومات الهوية المسرّبة.
\end{itemize}

\end{tcbwarn}

\textbf{التعامل الطارئ بعد الخطف — السرعة هي كل شيء}\par

\begin{enumerate}
  \item اتصل فورًا بخدمة العملاء للإبلاغ: China Mobile 10086 / China Unicom 10010 / China Telecom 10000، واطلب إيقاف الخدمة وتجميد استبدال البطاقة فورًا، لمنع المهاجم من متابعة استخدام رقمك.
  \item سجّل الدخول إلى الحسابات الحرجة من جهاز آخر: استخدم الكمبيوتر أو هاتف العائلة لتغيير كلمات مرور البنوك وAlipay وWeChat وApple ID فورًا، واطرد الأجهزة الغريبة.
  \item جمّد حسابات الدفع: Alipay 95188 للإبلاغ، ومركز أمان WeChat للتجميد، وخدمة العملاء المصرفية لتجميد الخدمات المصرفية عبر الإنترنت.
  \item اذهب إلى مكتب الخدمة ببطاقة الهوية: احضر بطاقتك الشخصية إلى المكتب الرسمي للحصول على بطاقة جديدة واستعادة حق استخدام الرقم.
  \item التبليغ للشرطة: خطف SIM يُعد «جريمة الاعتداء على المعلومات الشخصية للمواطنين»، لذا أحضر جميع الأدلة (سجلات وقت الإيقاف، الخسائر المالية، سجلات استبدال البطاقة) عند التبليغ.
  \item الفحص الشامل: بعد استعادة الرقم، افحص سجلات الدخول الأخيرة والمعاملات المالية واحدًا واحدًا، وتأكد من عدم وجود آثار استخدام غير مصرح به.
\end{enumerate}

\begin{tcbwarn}{الوقاية خير من العلاج — 5 أمور يومية}
\begin{enumerate}
  \item كلمة مرور شاشة القفل + التحقق الثنائي، فعّل كلاهما
  \item تأكد من تفعيل «البحث عن هاتفي»
  \item اضبط رمز PIN لبطاقة SIM (الإعدادات > الهاتف > رمز PIN لبطاقة SIM)، حتى لا يتمكن أحد من استخدام بطاقتك عند فقدان الهاتف
  \item لا تثبّت تطبيقات من مصادر غير معروفة، ولا تنقر على روابط غريبة
  \item خذ نسخًا احتياطية دورية للبيانات المهمة (الصور وجهات الاتصال وسجلات المحادثات)
\end{enumerate}

\end{tcbwarn}

\begin{tcbpractice}{تطبيق عملي: أعد «بطاقة الاتصال الطارئة» الخاصة بك}
\begin{enumerate}
  \item خذ ورقة واكتب الأرقام التالية: خدمة عملاء شركة الاتصالات (China Mobile 10086 / China Unicom 10010 / China Telecom 10000)، وخدمة العملاء المصرفية، و96110 لمكافحة الاحتيال، و110 للشرطة، وهواتف العائلة.
  \item اكتب بريد Apple ID/حساب العلامة التجارية المسجّل ووسائل الاتصال الاحتياطية.
  \item ضع هذه الورقة في المحفظة أو ألصقها على الثلاجة — ولا تحفظها في الهاتف، لأن الهاتف قد لا يكون معك عند الحادث.
  \item افحص الآن: هل «البحث عن هاتفي» مفعل؟ وهل رمز PIN لبطاقة SIM مضبوط؟
\end{enumerate}

\end{tcbpractice}

\begin{tcbsummary}{ملخص الفصل}
\begin{itemize}
  \item الـ 30 دقيقة الذهبية بعد فقدان الهاتف: الاتصال من هاتف آخر القفل عن بُعد إيقاف SIM تجميد الدفع إبلاغ الأهل التبليغ للشرطة
  \item الـ 72 ساعة بعد الاحتيال: إيقاف التحويل حفظ الأدلة الاتصال بـ 96110 التبليغ التواصل مع البنك تغيير كلمات المرور
  \item إصابة الهاتف: الوضع الآمن حذف التطبيقات المشبوهة الفحص الأمني إعادة ضبط المصنع (الحل النهائي)
  \item سرقة الحساب: تغيير كلمة المرور طرد الأجهزة فحص سجلات الدخول تفعيل التحقق الثنائي فحص الربط فحص الدفع
  \item رسائل الابتزاز 99.9\% احتيال، لا ترد ولا تدفع، التقط لقطة واحظر
  \item احتيال تبديل الوجه/تقليد الصوت: راقب الرمش/الشفاه/الحواف/الإضاءة، واطلب التحول وتغطية الوجه، واسأل أسئلة خاصة، واتصل بالرقم المعروف
  \item خطف SIM: عند اختفاء الإشارة فجأة كن يقظًا؛ اضبط PIN، وفعّل التحقق الثنائي، واستخدم تطبيقات المصادقة للحسابات المهمة
  \item الوقاية اليومية: شاشة القفل + 2FA، والبحث عن الهاتف، وPIN لبطاقة SIM، وعدم تثبيت تطبيقات مجهولة، والنسخ الاحتياطي الدوري
\end{itemize}

\end{tcbsummary}
